\SourceLeaf{281}%
\(x\). These coordinates are understood to be expressed in parts of the radius. They must therefore still be multiplied by \(A\) if they are wanted as lengths, or by \(206265''\) if they are wanted in angular measure.

An element of \(M\) is expressed by
\[
\BareRoot{(\cos y^2\mathrm{d}x^2+\mathrm{d}y^2)}
\]
or by \(\frac{\cos y}{\cos\psi}\mathord{.}\mathrm{d}x\) if we set
\[
\frac{\mathrm{d}y}{\cos y\mathrm{d}x}=\operatorname{tang}\psi
\]
so that \(\psi\) denotes the inclination of the element to the parallel to \(N\). To fix the picture, one may imagine the \(x\) increasing from right to left and the \(y\) from bottom to top. This determines the direction of positive \(\psi\).

The scale ratio for mapping the ellipsoid onto the sphere, denoted as above by \(m\), may here be regarded as a function of \(x\) and \(y\). The length of the element of \(L\) corresponding to that element of \(M\) will be
\[
=\frac{A\cos y}{m\cos\psi}\mathord{.}\mathrm{d}x
\]
and if, for brevity, we set
\[
\begin{aligned}
\log\operatorname{tang}(45^0+\tfrac{1}{2}y)&=u\\
\frac{\cos y}{m}&=n
\end{aligned}
\]
then \(n\) too will be a function of \(x\) and \(y\), or, equivalently, of \(x\) and \(u\). We then have
\[
\operatorname{tang}\psi=\frac{\mathrm{d}u}{\mathrm{d}x}
\]
and the element of \(L\)
\[
=\frac{An}{\cos\psi}\mathord{.}\mathrm{d}x
\]
\indent The nature of the line \(M\) is thus determined by the condition that, between any prescribed limits, the integral \(\int\frac{n}{\cos\psi}\mathrm{d}x\) or
\[
\int n\BareRoot{(1+\frac{\mathrm{d}u^2}{\mathrm{d}x^2})}\mathrm{d}x
\]
